\documentclass[../sablona-main.tex]{subfiles}

\begin{document}
	
	
	\section{Matematický model}
	
	\subsection{Model s konkrétními daty}
	Cílem je maximalizovat celkový zisk:
	\begin{align}
		\max\; Z &=\max\; (200x_1 + 350x_2).
	\end{align}
	
	Hlavní omezení (4 zdroje):
	
	\begin{enumerate}
		\item \textbf{Nohy:} každý stůl má 5 nohou, k dispozici je 300 nohou:
		\begin{align}
			5x_1 + 5x_2 \le 300.
		\end{align}
		
		\item \textbf{Dřevěné desky:} základní model používá 1 dřevěnou desku, k dispozici je 50:
		\begin{align}
			x_1 \le 50.
		\end{align}
		
		\item \textbf{Skleněné desky:} luxusní model používá 1 skleněnou desku, k dispozici je 35:
		\begin{align}
			x_2 \le 35.
		\end{align}
		
		\item \textbf{Montážní čas:} základní model $0{,}6$ h, luxusní $1{,}5$ h, k dispozici 63 h:
		\begin{align}
			0.6x_1 + 1.5x_2 \le 63.
		\end{align}
	\end{enumerate}
	
	Doména proměnných:
	\begin{align}
		x_1 \ge 0,\quad x_2 \ge 0, \quad x_1,x_2\in \mathbb{Z}.
	\end{align}
	
	\subsection{Tabulka dat}
	\begin{table}[h!]
		\centering
		\begin{tabular}{lccc}
			\toprule
			Zdroj / produkt & základní (1) & luxusní (2) & kapacita \\
			\midrule
			nohy & 5 & 5 & 300 \\
			dřevěné desky & 1 & 0 & 50 \\
			skleněné desky & 0 & 1 & 35 \\
			montáž (h) & 0.6 & 1.5 & 63 \\
			\midrule
			zisk (USD/ks) & 200 & 350 & -- \\
			\bottomrule
		\end{tabular}
		\caption{Vstupní data modelu.}
	\end{table}
	
	\subsection{Obecný tvar modelu}
	Obecný model pro rozdělení zdrojů mezi produkty za účelem maximalizace zisku lze zapsat jako:
	
	\begin{align}
		\max\; Z &=\max \sum_{j\in J} z_j x_j,\\
		\quad 
		\sum_{i\in I} a_{ij} x_j &\le b_i \qquad \forall i\in I,\\
		x_j &\ge 0 \qquad \forall j\in J .
	\end{align}
	
\end{document}